\subsection{Analyse von Kundenanforderungen}
\label{subsec:analyse-kundenanforderungen}

Die Analyse beginnt mit fachlichem Ziel, Nutzergruppen, Arbeitsabläufen und
Einsatzumgebung. Erst danach werden Anforderungen technisch eingeordnet.
Gemeinsame Datenbearbeitung kann etwa Backend und Persistenz erfordern, legt
aber noch nicht PostgreSQL fest.

Getrennt erfasst werden Browser- und Desktopzugang, Serverfunktionen,
Persistenz, Offline-Bedarf, native Funktionen, externe APIs, Betrieb,
Deployment, Sicherheit und Compliance. \texttt{desktop-local} enthält trotz
fehlendem Cloud-Backend weder einen lokalen Persistenzprovider noch ein lokales
Datenmodell oder eine Synchronisationsfunktion.
Tabelle~\ref{tab:customer-requirement-checklist} bündelt die Leitfragen.

\begin{table}[H]
  \centering
  \scriptsize
  \renewcommand{\arraystretch}{1.1}
  \caption{Kompakte Checkliste zur technikneutralen Anforderungsanalyse}
  \label{tab:customer-requirement-checklist}
  \begin{tabularx}{\textwidth}{@{}>{\raggedright\arraybackslash}p{0.18\textwidth}>{\raggedright\arraybackslash}p{0.43\textwidth}>{\raggedright\arraybackslash}X@{}}
    \toprule
    \textbf{Bereich} & \textbf{Leitfrage} & \textbf{Einfluss auf die Auswahl} \\
    \midrule
    Fachfunktion und Nutzung & Welche Abläufe, Nutzergruppen und Einsatzorte sind abzudecken?
      & Bestimmt Produktumfang; noch keine Profilfestlegung. \\
    Zugriffskanäle & Wird das Produkt im Browser, als Desktop-Anwendung oder über beide Kanäle genutzt?
      & Trennt Web-, Desktop- und Full-platform-Pfade. \\
    Zentrale Dienste & Müssen Funktionen oder gemeinsam genutzte Daten serverseitig bereitstehen?
      & Erfordert ein backendfähiges Cloudprofil. \\
    Persistenz und Offline-Betrieb & Werden zentrale relationale Daten benötigt; muss der Client ohne Netz arbeitsfähig sein?
      & PostgreSQL nur bei Backendbedarf; lokale Speicherung und Synchronisation sind separate Ergänzungen. \\
    Native Integration & Welche Betriebssystem-APIs, Geräte oder plattformspezifischen Oberflächen sind notwendig?
      & Bestimmt Tauri-Bedarf und Abstand zur nativen Baseline. \\
    Integration und Betrieb & Welche Fremdsysteme, Zielumgebungen, Skalierungs- und Betriebsanforderungen bestehen?
      & Erzeugt produktspezifische Adapter und Infrastrukturentscheidungen. \\
    Security und Compliance & Welche Schutz-, Nachweis-, Audit- oder regulatorischen Vorgaben gelten?
      & Erfordert zusätzliche Maßnahmen und Evidenz unabhängig vom Profil. \\
    \bottomrule
  \end{tabularx}
  \smallskip
  {\footnotesize Quelle: Eigene Darstellung.\par}
\end{table}


Das Ergebnis umfasst Plattformkomponenten und Abweichungen, etwa bei
Authentifizierung, Datenmodellen, Integrationen, Observability, Backups,
Infrastruktur oder Signing. Erst dann lässt sich zwischen passendem Profil,
begrenzter Erweiterung und einer Lösung außerhalb der Baseline entscheiden.
Die Methode folgt damit bewusst der Reihenfolge Anforderung, technische
Einordnung und erst anschließend Profil- oder Technologieauswahl.
\beginnerinput{workspace/statement/800_einsatz_transfer/830}
